
\subsubsection{6.1 Possible Explanations for SE Failure}

The observed failure of SE on Llama-3-8B-Base is consistent with the sampling degeneracy account: when 89.4\% of TriviaQA queries produce K=1, SE≈0 for the vast majority of queries. A query-level score distribution that is 89\% at or near zero has essentially no rank variation and cannot predict correctness. The anti-correlation on TriviaQA (AUROC=0.4735 < 0.5) is interpretable: among the 10.6\% of non-degenerate queries (K>1), queries that the model is genuinely uncertain about may be more likely to produce diverse outputs (K>1, higher SE), while high-confidence incorrect answers produce K=1 (lower SE). If genuine uncertainty weakly co-occurs with K>1 among the minority of diverse queries, SE becomes weakly anti-correlated with incorrectness in the combined distribution.

This account is correlational rather than causal. Establishing causality would require an experiment that independently varies sampling diversity—for example, by comparing base versus instruction-tuned variants, or sweeping sampling temperature—while holding other conditions fixed. Such experiments are not reported here and are identified as the primary future work.

The discrepancy between present results (SE AUROC 0.47–0.55) and Farquhar et al. (2024) (SE AUROC 0.72–0.79) is consistent with a model-type difference. Published SE evaluations included instruction-tuned Llama-2 variants and GPT-3; these models may produce higher sampling diversity due to RLHF-induced output variation, placing them in a regime where SE clustering produces K>1 for a larger proportion of queries. The exact instruction-tuning status and sampling diversity of the models used in Farquhar et al. (2024) is not confirmed in this study; the model-type explanation is a hypothesis rather than a verified fact.

KLE's sub-chance performance (AUROC=0.2642 on TriviaQA) is consistent with the rank-1 Laplacian explanation: near-identical samples produce near-zero eigenvalue sums, and systematic near-zero scores may invert rank ordering relative to true uncertainty. The specific mechanism producing this inversion is not fully characterized.

\subsubsection{6.2 Implications for UQ Benchmarking}

If the sampling degeneracy account is correct, it implies that the validity of SE and KLE as UQ estimators is contingent on the model producing semantically diverse outputs under the chosen sampling protocol. This diversity precondition was not explicitly reported or tested in the original SE paper or subsequent evaluations. The \texttt{degenerate_fraction} metric provides a computable diagnostic that can be reported alongside AUROC to characterize whether the diversity precondition is met.

The following practical guidance follows from the observed results, with the caveat that these recommendations are based on a single model (Llama-3-8B-Base) and two datasets at 500-sample scale:

\begin{itemize}
\item Token-probability achieves AUROC above 0.65 on both TriviaQA and NaturalQuestions and does not require sampling diversity.
\item SelfCheckGPT-NLI is competitive with token-probability on TriviaQA (AUROC=0.6862) but substantially underperforms on NaturalQuestions (AUROC=0.4508); its utility appears dataset-dependent.
\item SE and KLE perform at or below random chance on TriviaQA and only marginally above chance on NaturalQuestions when applied to Llama-3-8B-Base under N=10, temperature=1.0.
\end{itemize}

\subsubsection{6.3 Limitations}

\textbf{L1 — Scale coverage:} Only Llama-3-8B-Base is fully evaluated. The 70B experiment was not completed. Results may differ at 70B scale, and no cross-scale comparison is possible.

\textbf{L2 — Single model family:} Results are from one base model (Llama-3-8B). Generalization to other base models, other model families, or instruction-tuned models is not established.

\textbf{L3 — Sample size:} 500 samples per dataset (versus the planned 17,944 TriviaQA test and 3,610 NaturalQuestions validation). The reduced scale limits statistical precision and sub-group analysis. The NQ correctness rate of 19.4\% (97/500 correct) may widen confidence intervals due to class imbalance.

\textbf{L4 — Causal attribution:} The degenerate_fraction observation is a post-hoc diagnostic, not a preregistered variable. The causal claim that degeneracy causes SE failure requires controlled manipulation experiments not conducted in this study.

\textbf{L5 — Method coverage:} SEPs produced null results due to insufficient probe data; SelfCheckGPT-BERTScore produced constant scores. Two of six planned methods were effectively non-functional, leaving the benchmark incomplete.

\textbf{L6 — No instruct-model comparison:} The hypothesis that SE functions correctly on instruction-tuned Llama-3 models (due to higher sampling diversity) is not tested and remains a prediction.

